\chapter{The compile-time DSL}
\label{ch:dsl}

\maya{}'s most distinctive feature is its DSL --- a tiny embedded
language inside \cpp{} for describing UI trees. The DSL is unusual
because it lives entirely at the type level: an expression like

\begin{lstlisting}
t<"Hello"> | Bold | Fg<100, 180, 255>
\end{lstlisting}

is a \emph{single value} whose type encodes every choice you made
(\emph{the string}, \emph{bold}, \emph{the colour}). This chapter
shows how that works, why the compiler refuses certain expressions,
and how to escape into the runtime when you need to.

\begin{aside}
If you have written HTML or React, you have written a DSL: a small,
specialised vocabulary for the one thing you are doing (describing
UI). \maya{}'s DSL is the same idea, but it sits inside \cpp{}, runs
at compile time, and refuses to compile if you build an illegal
tree. The bottom-up explanation in this chapter is long, but the
payoff is being able to read any expression in any \maya{} program
and know what every token contributes to the type.
\end{aside}

\section{From scratch: what a DSL even is}
\label{sec:dsl:what}

The phrase ``Domain-Specific Language'' has two flavours in
programming.

An \textbf{external DSL} is a language with its own parser, lexer,
and often its own interpreter: SQL, regular expressions, JSON, CSS,
the shell. You write it in a different syntax from the host
language, and a separate engine reads it.

An \textbf{embedded} (or \emph{internal}) DSL is a sub-vocabulary
within the host language. It looks like normal code, but the
operators and identifiers have been overloaded to express a different
domain. JavaScript's \code{describe("thing", () => \dots)} from
Mocha is an embedded DSL for testing. \cpp{}'s \code{std::cout <<
thing} is the world's most ubiquitous embedded DSL --- the \code{<<}
operator does \emph{stream insertion} when applied to streams, not
bit-shift.

The \maya{} DSL is the embedded kind. Every piece of it is legal
\cpp{}; the magic is in operator overloads, variable templates, and
the fact that we have made the result types interesting enough that
the compiler can reason about them.

Four properties define a good embedded DSL:

\begin{enumerate}[leftmargin=*]
  \item \textbf{It reads like the thing it describes.}
    \code{v(t<"a">, t<"b">)} looks vertical because we picked
    \code{v} for that role.
  \item \textbf{It composes with the rest of the language.} You can
    pass a DSL fragment to a function, return one from a function,
    store one in a variable.
  \item \textbf{The host's type checker still works.} A mistake in
    the DSL is a host-language error, not a runtime ``DSL
    interpreter blew up''.
  \item \textbf{There is an escape hatch back to the host
    language.} \maya{}'s \code{dyn([]{...})} is exactly this.
\end{enumerate}

In the rest of the chapter we will see how \maya{} achieves all four.

\section{Non-type template parameters, properly}
\label{sec:dsl:nttp}

The foundation of the entire DSL is the \cpp{}20 / 26 ability to use
structured values as template parameters. We touched on this in
Chapter~\ref{ch:cpp}; here is the longer version, because the DSL
leans on it constantly.

A \textbf{template parameter} can be a type (\code{typename T}), but
it can also be a \emph{value}:

\begin{lstlisting}
template <int N> struct Array { int data[N]; };
Array<8>  a;   // N = 8
Array<16> b;   // N = 16, different type from a
\end{lstlisting}

The parameter \code{N} is an integer, not a type. The template is
instantiated once per distinct value of \code{N}, and each
instantiation is a separate type. \code{Array<8>} and \code{Array<16>}
do not interconvert.

\subsection{Why integers are not enough}

Integers, pointers, and enum values have always been allowed as
non-type template parameters (NTTPs). What you could \emph{not} do,
until \cpp{}20, was use a class type as an NTTP. There was no way to
write

\begin{lstlisting}
struct Point { int x; int y; };
template <Point P> struct Thing { /* ... */ };
Thing<Point{3, 4}> t;     // illegal in C++17
\end{lstlisting}

This matters for us because a UI description involves not just
numbers but \emph{compound data}: a string is a class type, a style
is a class type, a colour is three bytes. If we cannot put those in
template parameters, the DSL is severely limited.

\cpp{}20 lifted the restriction for \textbf{structural} class types:
any literal type whose members are all public, with no user-defined
oddities, can be an NTTP. \cpp{}26 extended it further. The
consequence is that \code{maya::dsl::Str}, \code{maya::dsl::CTStyle},
and \code{maya::dsl::BoxCfg} can all appear as template parameters,
and the DSL works.

\subsection{\code{Str}: the building block}

The simplest structural NTTP in \maya{} is \code{Str} ---
a fixed-size string class:

\begin{lstlisting}
template <std::size_t N>
struct Str {
    char data[N]{};
    consteval Str(const char (&s)[N]) {
        for (std::size_t i = 0; i < N; ++i) data[i] = s[i];
    }
    consteval std::size_t size() const { return N - 1; }
    constexpr operator std::string_view() const { return {data, N - 1}; }
};
\end{lstlisting}

Line by line:

\begin{itemize}[leftmargin=*]
  \item \code{template <std::size\_t N> struct Str}: parameterised
    by the length \code{N}.
  \item \code{char data[N]\{\}}: a fixed array of \code{N}
    characters, zero-initialised.
  \item \code{consteval Str(const char (\&s)[N])}: constructor that
    \emph{must} run at compile time, taking a reference to a
    \code{char} literal of length \code{N}. The compiler deduces
    \code{N} from the literal's size.
  \item The loop copies the literal into \code{data}.
\end{itemize}

Now you can write \code{Str{"Hello"}} as a constant expression. The
resulting value has a \emph{type} that depends on the length and a
\emph{value} that contains the bytes --- and crucially, both are
available to other templates.

\subsection{Putting \code{Str} into the type system}

\begin{lstlisting}
template <Str S>
struct TextNode {
    static constexpr std::string_view text() { return {S.data, S.size()}; }
};

TextNode<"Hello"> a;
TextNode<"World"> b;
// decltype(a) and decltype(b) are different types,
// even though both store no data per-instance.
\end{lstlisting}

\code{TextNode<"Hello">} and \code{TextNode<"World">} are distinct
types. The string contents are part of the \emph{type identity},
not a field. Two instances of \code{TextNode<"Hello">} occupy zero
bytes and are interchangeable.

Why does this matter? Because the compiler can now reason about your
strings: it can choose overloads based on which string was used, it
can refuse to mix incompatible strings, and it can fold the string's
length into compile-time arithmetic --- without you writing any of
that by hand.

\subsection{\code{CTStyle}: a structural style}

The second key NTTP type:

\begin{lstlisting}
struct CTStyle {
    bool has_fg = false, has_bg = false;
    uint8_t fg_r = 0, fg_g = 0, fg_b = 0;
    uint8_t bg_r = 0, bg_g = 0, bg_b = 0;
    bool bold_ = false, dim_ = false, italic_ = false;
    bool underline_ = false, strike_ = false, inverse_ = false;

    consteval CTStyle merge(CTStyle o) const { /* combine fields */ }
    [[nodiscard]] Style runtime() const       { /* materialise */ }
};
\end{lstlisting}

A \code{CTStyle} value is a \textbf{bag of compile-time style
flags}: bold yes/no, foreground colour yes/no plus three RGB bytes,
and so on. The struct is structural (all-public fields, no virtuals,
no custom comparison), so it qualifies as an NTTP. The
\code{merge} method is \code{consteval} --- it can be called inside
template arguments to compute new \code{CTStyle}s from old ones.

A style tag exists for every bool field:

\begin{lstlisting}
template <CTStyle V>
struct StyTag {
    static constexpr CTStyle value = V;
};

inline constexpr StyTag<CTStyle{.bold_ = true}>      Bold{};
inline constexpr StyTag<CTStyle{.dim_ = true}>       Dim{};
inline constexpr StyTag<CTStyle{.italic_ = true}>    Italic{};
inline constexpr StyTag<CTStyle{.underline_ = true}> Underline{};
\end{lstlisting}

\code{Bold} is a \code{StyTag} whose template argument is a
\code{CTStyle} with one field set. The variable \code{Bold} itself
is a zero-size instance --- the \code{CTStyle} lives in the type.

\subsection{\code{Fg<R,G,B>}: colours in the type system}

Colours are encoded the same way: a \code{CTStyle} with the
foreground fields set, exposed through a variable template:

\begin{lstlisting}
namespace detail_rgb {
    template <uint8_t R, uint8_t G, uint8_t B>
    consteval CTStyle make_fg() {
        return CTStyle{.has_fg=true, .fg_r=R, .fg_g=G, .fg_b=B};
    }
}

template <uint8_t R, uint8_t G, uint8_t B>
inline constexpr StyTag<detail_rgb::make_fg<R, G, B>()> Fg{};
\end{lstlisting}

\code{Fg<100, 180, 255>} is a zero-size value whose type is a
\code{StyTag} whose template parameter is a \code{CTStyle} whose
fields are exactly the RGB you asked for. The numbers fed in become
part of the type. Nothing about colour leaks into runtime data
until \code{.build()}.

\begin{funfact}
The MSVC compiler had (and partly still has) a long-standing bug
around \code{constexpr} evaluation of designated-initialised NTTPs.
If you look in \file{dsl.hpp} you will see a small namespace
\code{detail\_rgb} with \code{make\_fg} / \code{make\_bg} helpers ---
that existence-of-helpers is the workaround. The compiler bug ticket
is cited in a comment. Real libraries are full of small monuments
like this; each is a story of someone who hit the same wall before
you.
\end{funfact}

\begin{idea}
NTTPs let you put \emph{values} into the type system. \maya{} uses
that to encode strings (\code{Str}), styles (\code{CTStyle}),
colours (RGB triples), borders, paddings, and growth factors. Once a
value lives in a type, the compiler can check rules between values
at compile time --- which is the whole point.
\end{idea}

\section{The shape}

The DSL has six kinds of node:

\begin{center}
\small
\begin{tabular}{ll}
\code{TextNode<Str, CTStyle>}     & compile-time text \\
\code{RuntimeTextNode<S>}         & runtime text (from a string or number) \\
\code{BoxNode<Dir, Cfg, Cs...>}   & vertical or horizontal container \\
\code{DynNode<F>}                 & lambda escape hatch \\
\code{MapNode<R, P>}              & project a range to elements \\
\code{SpacerNode}                 & a flexible gap \\
\end{tabular}
\end{center}

Every node satisfies the \code{Node} concept:

\begin{lstlisting}
template <typename T>
concept Node = requires(const T& n) {
    { n.build() } -> std::convertible_to<Element>;
};
\end{lstlisting}

A \code{Node} is anything with a \code{.build()} that yields a runtime
\code{Element}. That is the only contract; the rest is how nodes
combine.

\section{\code{t<"Hello">} --- string in a type}

The simplest node is the compile-time text node:

\begin{lstlisting}
template <Str S, CTStyle Sty = CTStyle{}>
struct TextNode {
    static constexpr Str  str   = S;
    static constexpr CTStyle style = Sty;

    Element build() const {
        return TextElement{std::string{S.view()}, Sty.runtime()};
    }
    operator Element() const { return build(); }
};

template <Str S>
inline constexpr TextNode<S> t{};
\end{lstlisting}

The template parameters are: an \code{Str} (the string contents, an
NTTP from Chapter~\ref{ch:cpp}) and a \code{CTStyle} (a structural
type holding the compile-time style). \code{t<"Hello">} is the
variable template: \code{TextNode<"Hello", CTStyle\{\}>\{\}}.

Two facts to internalise:
\begin{itemize}[leftmargin=*]
  \item \code{t<"a">} and \code{t<"b">} have \emph{different types}.
    Their string contents live in the type system, not in a field.
  \item \code{TextNode} has no nonstatic data members. It is a zero-size
    type. Two instances of \code{TextNode<"a">} are
    indistinguishable; the compiler treats them as the same value.
\end{itemize}

\section{Pipes are operator overloads}

The DSL's pipe (\code{|}) is just \code{operator|}. \maya{} defines
overloads that match \code{(Node, StyTag)} and return a new node with
the style merged in:

\begin{lstlisting}
template <Str S, CTStyle Sty, CTStyle V>
constexpr auto operator|(TextNode<S, Sty>, StyTag<V>) {
    return TextNode<S, Sty.merge(V)>{};   // new type, same instance shape
}
\end{lstlisting}

Read the return type: the merged style is computed by \code{Sty.merge(V)},
which is \code{constexpr}, so the result is a brand-new
\code{TextNode<S, Sty'>}. Apply two style tags and you get yet another
type. The pipeline is sequential type rewriting at compile time.

\begin{idea}
Because the result of each pipe is a new type, the compiler can
enforce \emph{type-state} rules: a method is only available on a node
that has reached a particular state.
\end{idea}

This is the trick behind the border colour example. The pipe overload
for \code{BColTag} (a border colour tag) has a \code{requires} clause:

\begin{lstlisting}
template <FlexDirection Dir, BoxCfg Cfg, typename... Cs,
          uint8_t R, uint8_t G, uint8_t B>
    requires (Cfg.has_border)
constexpr auto operator|(BoxNode<Dir, Cfg, Cs...> n, BColTag<R, G, B>) {
    /* return a new BoxNode with border colour set */
}
\end{lstlisting}

If you call \code{v(...) | bcol<60, 65, 80>} without first applying
\code{border\_<Round>}, \code{Cfg.has\_border} is \code{false}, the
\code{requires} clause fails, no overload matches, and the compiler
prints an error. It is the same mechanism a Rust crate would use with
typestate traits.

\section{\code{v()} and \code{h()} --- variadic containers}

A box takes any number of children and a direction:

\begin{lstlisting}
template <FlexDirection Dir, BoxCfg Cfg, typename... Cs>
struct BoxNode {
    std::tuple<Cs...> children;

    Element build() const {
        BoxBuilder b{Dir, Cfg};
        std::apply([&](const auto&... cs) {
            (b.add(cs.build()), ...);   // fold-expression over children
        }, children);
        return b.finalize();
    }
};

template <typename... Cs>
constexpr auto v(Cs... cs) {
    return BoxNode<FlexDirection::Column, BoxCfg{}, Cs...>{
        std::tuple{std::move(cs)...}
    };
}

template <typename... Cs>
constexpr auto h(Cs... cs) {
    return BoxNode<FlexDirection::Row, BoxCfg{}, Cs...>{
        std::tuple{std::move(cs)...}
    };
}
\end{lstlisting}

Two \cpp{} features worth pausing on:

\begin{itemize}[leftmargin=*]
  \item \textbf{Variadic templates} (\code{typename... Cs}): you can
    declare ``zero or more types''. \code{Cs...} is a \emph{parameter
    pack}; you expand it with \code{...}.
  \item \textbf{Fold expressions} (\code{(b.add(cs.build()), ...)}):
    short for ``apply this expression to every element of the
    pack''. The comma is the \emph{fold operator} --- evaluate each,
    discard intermediate results.
\end{itemize}

The result: \code{v(t<"a">, t<"b">, t<"c">)} is a single
\code{BoxNode<Column, BoxCfg\{\}, TextNode<"a">, TextNode<"b">,
TextNode<"c">>}, holding a \code{tuple} of three zero-size text
nodes. Total runtime cost: zero bytes (the tuple of empty types is
empty).

\section{\code{BoxCfg} --- the type-state}

\code{BoxCfg} is a \code{constexpr} struct that records every layout
decision applied to a box:

\begin{lstlisting}
struct BoxCfg {
    int pad_top    = 0, pad_right = 0, pad_bottom = 0, pad_left = 0;
    int gap        = 0;
    int grow_num   = 0;
    bool has_border       = false;
    BorderStyle border    = BorderStyle::None;
    bool has_border_color = false;
    Color border_color{};
    Align align_items    = Align::Auto;
    Justify justify      = Justify::Auto;
    // ... etc.

    constexpr BoxCfg with_padding(int t, int r, int b, int l) const {
        BoxCfg c = *this;
        c.pad_top = t; c.pad_right = r; c.pad_bottom = b; c.pad_left = l;
        return c;
    }
    constexpr BoxCfg with_border(BorderStyle bs) const {
        BoxCfg c = *this; c.has_border = true; c.border = bs; return c;
    }
    // ...
};
\end{lstlisting}

Every layout pipe (\code{pad<1>}, \code{gap\_<1>}, \code{border\_<Round>},
\code{grow\_<2>}) returns a new \code{BoxNode} whose \code{BoxCfg}
template parameter has been replaced via the corresponding
\code{with\_*} method. The type carries the full configuration; the
final \code{.build()} reads the type and produces the runtime
\code{BoxElement}.

\section{Static, dynamic, and the bridge}

A pure-compile-time tree is fine for a banner but useless for a real
app: the count, the file list, the current time are all known only at
runtime. \maya{} provides three bridge nodes.

\subsection{\code{text(...)} --- runtime strings}

\begin{lstlisting}
text("hello")                // wraps a string_view
text(std::to_string(n))      // wraps a string
text(42)                     // wraps an int, formatted at build time
text(3.14)                   // wraps a double
\end{lstlisting}

\code{text} returns a \code{RuntimeTextNode<S>} --- a node whose
content is a runtime value but which still supports the DSL pipes.
Its \code{build()} formats the value and produces a
\code{TextElement}.

\subsection{\code{dyn([](){ ... })} --- arbitrary lambdas}

\begin{lstlisting}
dyn([&] {
    return loading ? text("Loading...") | Dim
                   : text("Ready")     | Bold | Fg<80, 220, 120>;
})
\end{lstlisting}

\code{dyn} wraps a lambda that returns anything convertible to
\code{Element}. \code{build()} just calls the lambda. This is the
universal escape hatch: when you cannot express something in the DSL
itself, write a lambda.

\subsection{\code{map(range, projection)} --- list comprehension}

\begin{lstlisting}
std::vector<File> files = ...;

auto list = v(
    t<"Files:"> | Bold,
    map(files, [](const File& f) {
        return h(text(f.name) | Bold, space, text(f.size) | Dim);
    })
);
\end{lstlisting}

\code{map} produces a \code{MapNode<R, P>}. At \code{build()} time it
walks the range, applies the projection, and produces an
\code{ElementList} that is spliced into the parent's children.

\section{Reading a real pipeline}

Put it all together. Here is a typical maya panel:

\begin{lstlisting}
auto panel = v(
    t<"Dashboard"> | Bold | Fg<100, 180, 255>,
    sep,
    dyn([&] { return text("Active users: " + std::to_string(state.users)); }),
    map(state.items, [](const auto& it) {
        return h(text(it.name) | Bold, space, text(it.value) | Dim);
    }),
    space,
    t<"[q] quit"> | Dim
) | border_<Round> | bcol<60, 65, 80> | pad<1, 2> | gap_<1>;
\end{lstlisting}

What is the type of \code{panel}? Approximately:

\begin{lstlisting}
BoxNode<
    FlexDirection::Column,
    BoxCfg{ /* round border, colour, pad 1/2, gap 1 */ },
    TextNode<"Dashboard", { bold, fg=(100,180,255) }>,
    SepNode,
    DynNode<lambda1>,
    MapNode<vector<Item>, lambda2>,
    SpacerNode,
    TextNode<"[q] quit", { dim }>
>
\end{lstlisting}

The compiler holds every choice you made --- the strings, the styles,
the border, the padding --- in this one large type. \code{panel.build()}
then walks it once, recursively, and produces an \code{Element} tree
that the layout engine and renderer take from there.

\section{Why compile-time?}

You could write this DSL with runtime polymorphism (each node a
\code{shared\_ptr<Node>}, each pipe a method that mutates fields). It
would work. But:

\begin{itemize}[leftmargin=*]
  \item \textbf{Allocations vanish.} A purely compile-time tree of
    zero-size types has no heap traffic.
  \item \textbf{Errors are at compile time.} Border colour without a
    border, negative padding (template parameter constraint
    \code{N >= 0}), wrong type of child --- all rejected before the
    program runs.
  \item \textbf{The compiler can fold.} A \code{constexpr} tree can
    be assigned to a \code{constexpr} variable and rendered at startup
    with no parsing.
\end{itemize}

The trade-off: error messages can be long. The first time you see a
500-line template error, breathe. The bottom of the message --- the
\code{required from here} line and the constraint that failed --- is
usually the actionable bit.

\section{Where the runtime takes over}

\code{.build()} is the boundary. Above it: types, NTTPs, fold
expressions, the world the compiler sees. Below it: \code{Element}, a
runtime value the layout and renderer understand. The next chapter
goes through that runtime side.

\begin{recap}
\maya{}'s DSL is a set of node types that compose with \code{|} into
larger node types. Strings, styles, borders, paddings, and tree
structure all live as template parameters. Type-state constraints
make invalid trees uncompilable. Runtime content enters through
\code{text}, \code{dyn}, and \code{map}. \code{.build()} converts the
compile-time tree into a runtime \code{Element}.
\end{recap}

\section{Every pipe, named and explained}
\label{sec:dsl:pipes-tour}

The DSL has roughly twenty overloads of \code{operator|}. Reading
the source for each takes ten seconds; reading them together is
useful, because the table makes it obvious which kinds of node
accept which kinds of modifier. Here is the full list, grouped by
the left-hand side.

\subsection{Pipes that take a \code{TextNode}}

\begin{center}
\small
\begin{tabular}{lll}
\textbf{Pipe} & \textbf{Effect} & \textbf{Result type} \\
\hline
\code{TextNode | StyTag<V>}    & merge style V       & \code{TextNode<S, Sty.merge(V)>} \\
\code{TextNode | WidthTag<W>}  & wrap in fixed box   & \code{BoxNode<Row, ct\_width=W>} \\
\code{TextNode | HeightTag<H>} & wrap in fixed box   & \code{BoxNode<Row, ct\_height=H>} \\
\end{tabular}
\end{center}

Text nodes accept any style modifier and nothing else. Adding
\code{pad<1>} to a bare text node is a compile error --- padding is
a container property; text has no notion of inset space.

\subsection{Pipes that take a \code{RuntimeTextNode}}

The runtime text node (built by \code{text(...)}) gets the same
style pipes plus two wrap-mode pipes:

\begin{center}
\small
\begin{tabular}{ll}
\textbf{Pipe} & \textbf{Effect} \\
\hline
\code{| StyTag<V>}    & merge style into the runtime node \\
\code{| clip}         & truncate with ellipsis on overflow \\
\code{| nowrap}       & never break the line, may overflow \\
\end{tabular}
\end{center}

The wrap-mode pipes do not exist on compile-time \code{TextNode}
because compile-time text is meant to be short and known; the
ellipsis case only matters when you wrap a runtime value of unknown
length.

\subsection{Pipes that take a \code{BoxNode}}

This is where the real type-state machine lives. A box accepts every
kind of modifier:

\begin{center}
\small
\begin{tabular}{lll}
\textbf{Pipe} & \textbf{Effect} & \textbf{Notes} \\
\hline
\code{| StyTag<V>}        & merge style on the box itself      & all boxes \\
\code{| PadTag<T,R,B,L>}  & set inner padding                  & all boxes \\
\code{| GapTag<G>}        & set gap between children           & all boxes \\
\code{| BorderTag<BS>}    & set border style                   & enables \code{bcol} pipe \\
\code{| BColTag<R,G,B>}   & set border colour                  & \textbf{requires} \code{has\_border} \\
\code{| GrowTag<G>}       & set flex grow factor               & all boxes \\
\code{| WidthTag<W>}      & set fixed width                    & all boxes \\
\code{| HeightTag<H>}     & set fixed height                   & all boxes \\
\end{tabular}
\end{center}

Four of the modifier types have a \code{static\_assert} on the
template parameter to reject obvious nonsense at compile time:

\begin{itemize}[leftmargin=*]
  \item \code{PadTag}: every value must be \code{>= 0}. Negative
    padding is a category error.
  \item \code{GapTag}: \code{G >= 0}.
  \item \code{GrowTag}: \code{G >= 0}.
  \item \code{WidthTag}, \code{HeightTag}: \code{> 0}. A width of
    zero is a hidden box, which we have a different way to express
    (\code{visible(false)}).
\end{itemize}

The one with the named requires-clause is \code{BColTag}:

\begin{lstlisting}
template <FlexDirection Dir, BoxCfg Cfg, typename... Cs,
          uint8_t R, uint8_t G, uint8_t B>
    requires (Cfg.has_border)
constexpr auto operator|(BoxNode<Dir, Cfg, Cs...> n, BColTag<R, G, B>);
\end{lstlisting}

The constraint \code{requires (Cfg.has\_border)} forbids the
overload when the box has no border yet. We will look at the error
message in section~\ref{sec:dsl:errors}.

\subsection{The style-merging pipe}

A quietly powerful one: two style tags combine into one.

\begin{lstlisting}
template <CTStyle A, CTStyle B>
constexpr auto operator|(StyTag<A>, StyTag<B>) {
    return StyTag<A.merge(B)>{};
}
\end{lstlisting}

This is the foundation of building reusable style presets:

\begin{lstlisting}
constexpr auto heading = Bold | fg<0xFFFFFF>;
constexpr auto error   = Bold | fg<0xFF4444>;
constexpr auto subtle  = Dim  | Italic;

t<"Title">     | heading
text(msg)      | error
t<"caption">   | subtle
\end{lstlisting}

The combined tag is itself a \code{StyTag<...>}, so it composes
further: \code{heading | Underline} is legal. Compose your design
system this way and applying it is a single token.

\subsection{The hex-colour shorthand}

\code{Fg<R, G, B>} is verbose for designers, who think in hex. The
DSL provides \code{fg<Hex>} and \code{bg<Hex>}:

\begin{lstlisting}
t<"warn">  | fg<0xFFAA00>
t<"error"> | fg<0xFF4444>
v(...)     | bg<0x1A1A2E>
\end{lstlisting}

The \code{consteval} helper that decodes the hex into RGB also
validates it: \code{fg<0xFFFFFFFF>} (32 bits) throws a string
literal from inside \code{consteval}, which the compiler reports as
a constant-expression failure --- ``\code{fg<>: hex value exceeds
0xFFFFFF}''. Compile-time error, not a silent black box.

\section{The type-state machine, drawn out}
\label{sec:dsl:state-machine}

The phrase \textbf{type-state} comes from research on programming
languages: encode the legal sequence of operations on an object into
its type, so that calling a method at the wrong stage is a type
error. Rust's \code{Builder} idioms are typical: the constructor
returns a \code{Builder<Empty>}, the \code{name(\dots)} method
returns a \code{Builder<Named>}, and \code{build()} is only defined
on \code{Builder<Named>}.

\maya{}'s \code{BoxNode} works the same way, with \code{BoxCfg} as
the state.

\paragraph{Initial state.} \code{v(\dots)} returns a
\code{BoxNode<Column, BoxCfg\{\}, \dots>}: every flag in
\code{BoxCfg} is false / zero. From this state:

\begin{itemize}[leftmargin=*]
  \item Padding, gap, grow, width, height, raw style: all available.
    Each returns a new \code{BoxNode} whose \code{BoxCfg} has the
    corresponding field updated.
  \item Border: available. Returns a \code{BoxNode} with
    \code{has\_border = true}.
  \item Border colour: \emph{not} available; the requires-clause
    fails because \code{has\_border} is false.
\end{itemize}

\paragraph{After \code{| border\_<Round>}.} \code{has\_border}
becomes true. Border colour becomes available.

\paragraph{Drawn as a state diagram.}

\begin{center}
\small
\begin{tabular}{rcl}
 \code{BoxNode<Cfg(has\_border=false)>}
  & $\xrightarrow{\;\;|\;\code{border\_<R>}\;\;}$
  & \code{BoxNode<Cfg(has\_border=true)>} \\
 \code{BoxNode<Cfg(has\_border=false)>}
  & $\xrightarrow{\;\;|\;\code{bcol<\dots>}\;\;}$
  & \textsc{compile error} \\
 \code{BoxNode<Cfg(has\_border=true)>}
  & $\xrightarrow{\;\;|\;\code{bcol<\dots>}\;\;}$
  & \code{BoxNode<\dots, has\_border\_color=true>} \\
\end{tabular}
\end{center}

The state machine is tiny --- two states for the border flag --- but
the pattern scales. Any rule of the form ``\emph{X} is only legal
after \emph{Y}'' can be encoded the same way, by toggling a field in
the \code{BoxCfg} on the \emph{Y} operation and adding a
\code{requires (Cfg.\emph{field})} to the \emph{X} overload.

\begin{idea}
Type-state in a DSL means: the legal sequence of pipes is encoded
into the chain of types. The compiler refuses pipes the current type
does not permit. The error has a name (a \code{requires} clause) and
a single source line.
\end{idea}

\section{Decoding the compile errors}
\label{sec:dsl:errors}

A powerful type system gets you powerful errors. The first time you
see a 200-line template error, the impulse is to scroll up and
panic; the trick is knowing where to look.

The shape of every constraint-failure error from the DSL is:

\begin{enumerate}[leftmargin=*]
  \item A line saying \emph{no viable overload of \code{operator|}}.
  \item One or more notes saying \emph{candidate template ignored:
    constraints not satisfied}.
  \item For each candidate, the constraint that did not hold:
    \code{Cfg.has\_border} was \code{false}, \code{T >= 0} did not
    hold for \code{T = -1}, and so on.
  \item Finally, the line in \emph{your} code that triggered it.
\end{enumerate}

\subsection{The border-colour error}

The source mistake:

\begin{lstlisting}
auto ui = v(t<"hi">) | bcol<60, 65, 80>;     // forgot border first
\end{lstlisting}

The heart of the error from a modern \code{g++}:

\begin{lstlisting}[style=shell]
error: no match for 'operator|' (operand types are
    'maya::dsl::BoxNode<maya::FlexDirection::Column,
      maya::dsl::BoxCfg{...has_border=false...}, ...>'
    and 'const maya::dsl::BColTag<60, 65, 80>')
note: candidate: 'template<...> constexpr auto
    maya::dsl::operator|(BoxNode<Dir, Cfg, Cs...>, BColTag<R, G, B>)
        requires (Cfg.has_border)'
note: constraints not satisfied
note:   the expression 'Cfg.has_border' evaluated to 'false'
\end{lstlisting}

The relevant fact is in the last note: \code{Cfg.has\_border}
evaluated to false. Fix:

\begin{lstlisting}
auto ui = v(t<"hi">) | border_<Round> | bcol<60, 65, 80>;
\end{lstlisting}

\subsection{The negative-padding error}

The source mistake:

\begin{lstlisting}
auto ui = v(t<"hi">) | pad<-1>;
\end{lstlisting}

The error:

\begin{lstlisting}[style=shell]
error: static assertion failed: Padding values must be non-negative
  in instantiation of 'struct maya::dsl::PadTag<-1, -1, -1, -1>'
\end{lstlisting}

A \code{static\_assert} caught it before any pipe overload was
even considered. The error names exactly what is wrong.

\subsection{The wrong-child-type error}

The DSL's \code{v} and \code{h} take children that satisfy the
\code{DslChild} concept --- either a \code{Node} (anything with
\code{.build()}) or a \code{std::ranges::range} of \code{Element}.
Passing something else fails the concept check:

\begin{lstlisting}
auto bad = v(42);                       // int is not a DslChild
\end{lstlisting}

The error roughly:

\begin{lstlisting}[style=shell]
error: no matching function for call to 'v(int)'
note: candidate template ignored: constraints not satisfied [with Cs = {int}]
note: because 'int' does not satisfy 'DslChild'
\end{lstlisting}

The fix is to wrap the value in \code{text(\dots)} or its appropriate
constructor: \code{v(text(42))}.

\subsection{Strategy for reading a template error}

Three rules that hold for every \maya{} compile failure:

\begin{enumerate}[leftmargin=*]
  \item \textbf{Read the bottom first.} The line that says
    \emph{required from here} or \emph{in the expansion of \dots}
    is closer to your source than the wall of template-instantiation
    backtrace above it.
  \item \textbf{Look for ``constraints not satisfied'' / ``static
    assertion failed''.} Those are the human-readable handles. The
    text after them names the rule that broke.
  \item \textbf{Identify the offending NTTP value.} The compiler
    will print the \code{BoxCfg} or \code{CTStyle} with its fields;
    the field whose value is wrong (\code{has\_border=false},
    \code{pad\_t=-1}) is the actionable bit.
\end{enumerate}

Most of the wall of text is irrelevant; you are looking for two or
three lines.

\begin{warning}
Do \emph{not} \code{\#include <maya/maya.hpp>} into a translation
unit that you then refuse to use because the error is intimidating.
The errors are intimidating because they are precise. The fix is
usually one token away; scrolling past the error makes finding it
harder.
\end{warning}

\section{What \code{.build()} actually does}
\label{sec:dsl:build}

The boundary between compile time and runtime is the \code{.build()}
method. Every node provides it; the implementation depends on the
kind of node.

\paragraph{\code{TextNode<S, Sty>::build()}.}

\begin{lstlisting}
Element build() const {
    return Element{TextElement{
        .content = std::string{std::string_view{S}},
        .style   = Sty.runtime(),
    }};
}
\end{lstlisting}

The compile-time string \code{S} and style \code{Sty} are read out
of the type. \code{S} becomes a \code{std::string} (heap-allocated
if the small-string optimisation does not apply; for short banners
it usually does). \code{Sty.runtime()} converts the
\code{CTStyle} bag into a runtime \code{Style} object (Chapter~\ref{ch:elements}).
The result is a \code{TextElement} wrapped in the \code{Element}
variant.

\paragraph{\code{RuntimeTextNode<S>::build()}.}

\begin{lstlisting}
Element build() const {
    if constexpr (std::integral<S> && !std::same_as<S, bool>) {
        // format integer via std::to_chars
    } else if constexpr (std::floating_point<S>) {
        // format float via std::to_chars, 2 fractional digits
    } else {
        // S is a string-like: copy into a std::string
    }
}
\end{lstlisting}

The \code{if constexpr} picks the right path at compile time, so
there is no per-call branching. \code{text(42)} formats with no
heap allocation (uses a 32-byte stack buffer); \code{text("hi")}
allocates a \code{std::string}.

\paragraph{\code{BoxNode<Dir, Cfg, Cs...>::build()}.} The most
involved:

\begin{lstlisting}
Element build() const {
    return std::apply([](const auto&... cs) {
        auto b = maya::detail::box().direction(Dir);
        if constexpr (Cfg.pad_t || Cfg.pad_r || ...)
            b = std::move(b).padding(Cfg.pad_t, ...);
        if constexpr (Cfg.gap > 0)
            b = std::move(b).gap(Cfg.gap);
        if constexpr (Cfg.has_border)
            b = std::move(b).border(Cfg.bstyle);
        if constexpr (Cfg.has_border_color)
            b = std::move(b).border_color(Color::rgb(Cfg.bc_r, Cfg.bc_g, Cfg.bc_b));
        // ... style, dimensions, etc.
        return b(cs.build()...);   // fold-expression: build each child
    }, children);
}
\end{lstlisting}

Notice the \code{if constexpr} on every \code{BoxCfg} field. A box
with no border generates \emph{no} call to \code{.border()}; the
branch does not exist in the resulting binary. The compiler
specialises one \code{build()} per distinct \code{BoxCfg} and
throws out everything irrelevant.

The trailing \code{b(cs.build()...)} expands the children pack:
every child's \code{.build()} is invoked, and the resulting
\code{Element}s are handed to the runtime box builder via its
function-call operator (Chapter~\ref{ch:elements}).

\paragraph{\code{DynNode<F>::build()}.}

\begin{lstlisting}
Element build() const { return fn(); }
\end{lstlisting}

Literal one-liner. The lambda is invoked; whatever it returns is
converted to \code{Element}.

\paragraph{\code{MapNode<R, Proj>::build()}.}

\begin{lstlisting}
Element build() const {
    std::vector<Element> items;
    if constexpr (std::ranges::sized_range<R>) items.reserve(range.size());
    for (auto&& val : range) items.emplace_back(proj(val).build());
    return Element{ElementList{std::move(items)}};
}
\end{lstlisting}

The range is walked, the projection applied, and the results
gathered into an \code{ElementList}. Because \code{MapNode} produces
an \code{ElementList} (not a single \code{Element}), the parent
\code{v}/\code{h} splices its children into the parent's child list
--- the moral equivalent of a JSX fragment.

\subsection{The cost of \code{.build()}}

A pure compile-time tree (\code{t<"\dots">} and \code{v}/\code{h}
with no runtime children) does \code{.build()} in:

\begin{itemize}[leftmargin=*]
  \item One \code{std::string} construction per text node
    (small-string optimisation often makes this free).
  \item One \code{BoxElement} construction per box node, which holds
    a \code{std::vector<Element>} of children. The vector is
    reserved to the right size.
  \item Some \code{Style} field assignments.
\end{itemize}

No virtual dispatch, no template-method indirection. The whole tree
is built in a single recursive sweep that the compiler often inlines
completely for small trees. For a typical panel with twenty nodes,
\code{.build()} runs in tens of microseconds.

The \code{Element} tree is then dropped on the floor at the end of
the frame (after rasterisation into the canvas), and the next frame
builds a fresh one. This is the rebuild-every-frame discipline from
Chapter~\ref{ch:elm}.

\begin{recap}
\code{.build()} is the bridge from the compile-time tree to the
runtime \code{Element}. It is a one-shot recursive walk: every
node's \code{build()} produces an \code{Element}, children are
gathered into the parent's \code{Element}, and a runtime tree comes
out the bottom. Then it is the runtime side's problem.
\end{recap}

\section{Recipe gallery}
\label{sec:dsl:recipes}

A collection of patterns you will use constantly. Each is a small,
composable fragment; the goal is to make them feel familiar before
you need them.

\subsection{A titled panel with a border}

\begin{lstlisting}
constexpr auto title_panel = v(
    t<"Status"> | Bold | Underline,
    sep,
    t<"all systems nominal"> | Fg<80, 220, 120>
) | pad<1, 2> | border_<Round> | bcol<60, 65, 80>;
\end{lstlisting}

\code{sep} is a horizontal separator (a one-row box with a top
border); \code{pad<1, 2>} is one row of vertical padding, two
columns of horizontal.

\subsection{A two-column layout that grows to fit}

\begin{lstlisting}
auto split = h(
    sidebar  | w_<24>,
    content  | grow_<1>
) | grow_<1>;
\end{lstlisting}

The sidebar has fixed width 24; the content gets every remaining
column. The outer box grows to fill whatever vertical container is
holding it.

\subsection{A list of items from a runtime collection}

\begin{lstlisting}
auto list = v(
    t<"Files"> | Bold,
    map(files, [](const File& f) {
        return h(text(f.name) | w_<30>,
                 text(f.size) | Dim);
    })
);
\end{lstlisting}

The lambda is called once per file; each call produces an
\code{Element}. \code{map} gathers them into an \code{ElementList}
that \code{v} splices in.

\subsection{A conditional region}

Two forms, both common:

\begin{lstlisting}
// when(cond, then) --- hidden if cond is false
auto top = v(
    title,
    when(show_advanced, advanced_panel)
);

// when(cond, then, else) --- one or the other
auto status = when(loading,
    text("loading\dots") | Dim,
    text("ready")        | Fg<80, 220, 120>);
\end{lstlisting}

The one-argument form falls back to \code{BlankNode\{\}} (a
zero-cell empty text). The two-argument form picks one of two
branches; both branches must satisfy \code{Node}.

\subsection{A reusable style preset}

\begin{lstlisting}
constexpr auto heading = Bold | fg<0xFFFFFF>;
constexpr auto error   = Bold | fg<0xFF4444>;
constexpr auto subtle  = Dim  | Italic;

auto ui = v(
    t<"Welcome"> | heading,
    text(err)    | error,
    t<"v0.1">    | subtle
);
\end{lstlisting}

The presets are themselves \code{StyTag}s, computed at compile time.
Applying one is a single token; defining your design system this
way is the path to consistency.

\subsection{Inline width hints on text}

\begin{lstlisting}
text(very_long_string) | clip | w_<30>
\end{lstlisting}

\code{clip} switches the text node into truncate-with-ellipsis
mode. \code{w\_<30>} wraps the result in a fixed-width box, so the
text is forced to 30 columns regardless of available space.

\subsection{An animated spinner}

\begin{lstlisting}
dyn([&] {
    constexpr std::array frames = {"top", "upper-right", "lower-right", "lower-left"};
    return text(frames[(tick / 100) % frames.size()]);
}) | Fg<140, 220, 255>
\end{lstlisting}

The lambda picks a frame from the array based on the current tick
count. Combine with \code{Sub<Msg>::on\_animation\_frame(Tick\{\})}
from Chapter~\ref{ch:elm} for a 60\,Hz spinner.

\subsection{A right-aligned status bar}

\begin{lstlisting}
auto status_bar = h(
    text(left_text),
    space,                          // SpacerNode: grows to fill
    text(right_text) | Dim
);
\end{lstlisting}

\code{space} is the spacer node: a zero-content box with
\code{grow=1}. Between two non-growing siblings it stretches to
push them to opposite ends of the row.

\section{Why all of this, again?}
\label{sec:dsl:why-again}

At the end of a long chapter on templates, fold expressions, and
type-state machines, it is worth stepping back. What did all that
machinery buy us?

At build time:

\begin{itemize}[leftmargin=*]
  \item \textbf{No allocation for the tree shape.} The structure of
    your UI lives in types; only the leaf values (runtime strings,
    runtime children) allocate.
  \item \textbf{Errors are catch-able before the program runs.}
    Border colour without border, negative padding, wrong child
    type, hex overflow --- all rejected by the compiler.
  \item \textbf{The DSL composes with the host language.} Helper
    functions, lambdas, ranges --- all work without special
    integration.
\end{itemize}

At runtime:

\begin{itemize}[leftmargin=*]
  \item \textbf{\code{.build()} is one linear walk}, specialised by
    the compiler per distinct configuration. No virtual dispatch.
  \item \textbf{The runtime side is simple.} It receives a plain
    tree of \code{Element} values and lays it out; it does not
    care that the tree was generated by a compile-time DSL.
\end{itemize}

For build time. For correctness. For speed. Three things at once,
because we paid the cost in the DSL design.

\begin{idea}
The \maya{} DSL is what you get when you take ``the compiler should
catch this'' seriously and follow it all the way down. Strings,
styles, sizes, structure: in the types. Mistakes: caught at build.
Runtime: cheap and uniform. The trade is some up-front complexity in
the DSL itself; the payoff is every \maya{} program you ever write.
\end{idea}

\section*{Workshop}

\begin{enumerate}[leftmargin=*]
  \item Write a panel that shows your name in bold, an underlined
    tagline, and a divided list of three hobbies. Use only the
    compile-time pipes.
  \item Add a runtime counter (an \code{int} variable) with
    \code{dyn}. Increment it in a loop with
    \code{maya::live(...)} (covered in Chapter~\ref{ch:cmdsub}) and
    watch the panel refresh.
  \item Try \code{bcol<60, 65, 80>} on a box without a border. Read
    the compiler error. Find the line that says \code{constraints
    not satisfied} and the \code{Cfg.has\_border} that triggered it.
\end{enumerate}
